\documentclass[../../main.tex]{subfiles}
\begin{document}
\section{带角流形}\label{section:ch16-manifolds-with-corners}

为了把 Stokes 定理应用于像立方体、单形这样的区域，我们需要比带边流形稍稍一般一点的一类空间。令 $\mathbb{H}^n_+$ 表示 $\mathbb{R}^n$ 的闭第一正卦限：
\begin{align*}
\mathbb{H}^n_+=\{(x^1,\dots,x^n)\in\mathbb{R}^n:x^1\geqslant0,\dots,x^n\geqslant0\}.
\end{align*}
$\mathbb{H}^n_+$ 中的一个点称为\emph{边界点}，如果它至少有一个坐标为零；其余点称为\emph{内点}。

\begin{definition}[带角流形]\index{流形!带角}
一个 $n$ 维\emph{带角光滑流形}是一个第二可数的 Hausdorff 空间，它在每一点处都局部微分同胚于 $\mathbb{H}^n_+$ 的一个开子集。从 $\mathbb{H}^n_+$ 的开子集到 $\mathbb{H}^n_+$ 的开子集的光滑映射，是指它可以光滑地延拓到 $\mathbb{R}^n$ 的某个开子集上。
\end{definition}

每个带边流形自动地是带角流形（因为 $\mathbb{H}^n$ 可以嵌入 $\mathbb{H}^n_+$ 作为一个面）。反过来一般不成立：例如，立方体 $[0,1]^n$ 和标准单形 $\Delta^n$ 都是带角流形，但不是带边流形。

在一个带角流形 $M$ 上，一个点被称作是一个 $k$ 角点，如果在某个（因而每个）角坐标卡中它恰有 $k$ 个坐标为零。$0$ 角点就是内点，$1$ 角点是普通边界点，$2$ 角点及以上是真正的角点。所有 $k$ 角点的集合记为 $\partial_k M$。下面的命题说明这个概念是内蕴的。

\begin{proposition}[角点的不变性]\label{proposition:ch16-invariance-of-corner-points}
设 $M$ 是一个带角光滑流形。一个点是 $k$ 角点这一性质与角坐标卡的选取无关。
\end{proposition}

\begin{proposition}\label{proposition:ch16-corner-parametrization}
设 $M$ 是一个定向 $n$ 维带角光滑流形。$\partial_1 M$（一维角点以外的普通边界点）是一个 $(n-1)$ 维带边光滑流形，并带有诱导定向。每个 $k$ 角点的集合 $\partial_k M$ 是一个 $(n-k)$ 维带角光滑流形。
\end{proposition}

\begin{exercise}\label{exercise:ch16-22}
证明：一个 $n$ 维带角流形的边界 $\partial M=\partial_1 M\cup\partial_2 M\cup\cdots\cup\partial_n M$ 是一个 $(n-1)$ 维带角光滑流形。
\end{exercise}

\begin{example}[曲线段的参数化]\label{example:ch16-curve-segments-param}
设 $I=[0,1]$。分段光滑曲线段 $\gamma:[a,b]\to M$ 的边界点集 $\partial I=\{0,1\}$。在例~16.23 中我们用到过如下约定：当用一个区间参数化一条有向曲线时，$\partial I$ 的两个端点带有符号：起点 $0$ 带 $-1$，终点 $1$ 带 $+1$。于是 $\int_{\partial I}\varphi=\varphi(1)-\varphi(0)$。
\end{example}

\begin{figure}[ht]
\centering
\includegraphics[width=0.6\textwidth]{fig-16-10.png}
\caption{带角流形的边界坐标卡}\label{figure:ch16-corner-chart}
\end{figure}

\begin{figure}[ht]
\centering
\includegraphics[width=0.6\textwidth]{fig-16-11.png}
\caption{立方体的边界定向}\label{figure:ch16-cube-boundary}
\end{figure}

Stokes 定理的证明几乎原封不动地推广到带角流形，因为唯一用到边界坐标卡的地方是在 $x^n=0$ 的那个面上作积分，而角点只出现在低维集合上，对积分没有贡献。

\begin{theorem}[带角流形上的 Stokes 定理]\label{theorem:ch16-stokes-with-corners}
设 $M$ 是一个定向 $n$ 维带角光滑流形，$\omega$ 是 $M$ 上一个紧支集光滑 $(n-1)$-形式。则
\begin{align*}
\int_M\mathrm{d}\omega=\int_{\partial M}\omega.
\end{align*}
\end{theorem}

\begin{figure}[ht]
\centering
\includegraphics[width=0.6\textwidth]{fig-16-12.png}
\caption{单形的边界}\label{figure:ch16-simplex-boundary}
\end{figure}

\begin{figure}[ht]
\centering
\includegraphics[width=0.6\textwidth]{fig-16-13.png}
\caption{带角坐标卡上的 Stokes 定理}\label{figure:ch16-stokes-corner-chart}
\end{figure}

\begin{theorem}\label{theorem:ch16-stokes-manifold-corners}
设 $M$ 是一个定向光滑流形（不带边），$D\subseteq M$ 是一个闭的、具有分段光滑边界的区域。给 $\partial D$ 以诱导定向。如果 $\omega$ 是 $M$ 上一个紧支集光滑 $(n-1)$-形式，则
\begin{align*}
\int_D\mathrm{d}\omega=\int_{\partial D}\omega.
\end{align*}
\end{theorem}

\begin{corollary}\label{corollary:ch16-simply-connected-exact}
在一个单连通光滑流形上，每个闭 $1$-形式都是恰当的。
\end{corollary}

\begin{proof}
设 $M$ 单连通，$\omega$ 是 $M$ 上的闭 $1$-形式。由于 $M$ 上每条分段光滑闭曲线都与一条常值曲线道路同伦，上面的定理表明 $\omega$ 在每条这样的曲线上的积分都等于 $0$。因此 $\omega$ 是保守的，从而是恰当的。
\end{proof}
\end{document}
